\documentclass{amsart}
\usepackage{amsmath,amssymb,amsthm,hyperref}
\usepackage{algorithm}
\usepackage{algpseudocode}

\theoremstyle{plain}
\newtheorem{theorem}{Theorem}
\newtheorem{lemma}{Lemma}
\newtheorem{proposition}{Proposition}
\newtheorem{corollary}{Corollary}
\theoremstyle{definition}
\newtheorem{definition}{Definition}
\newtheorem{remark}{Remark}

\newcommand{\dG}{d_{G'}}
\newcommand{\dH}{d_{H}}

\begin{document}

\title{Enumeration of geodetic homeomorphs of the Hoffman--Singleton graph}
\maketitle

\section{Statement and result}

Let $H$ denote the Hoffman--Singleton graph: the unique $(7,5)$--Moore graph,
which is $7$--regular on $50$ vertices, has $175$ edges, girth $5$, and diameter
$2$. For a positive integer assignment $x=(x_e)_{e\in E(H)}\in\mathbb{Z}_{\ge1}^{E(H)}$,
let $H^{x}$ be the graph obtained from $H$ by replacing each edge $e=uv$ with an
internally disjoint path of length $x_e$ joining $u$ and $v$. The vertices of $H$
are the \emph{branch vertices} of $H^x$; the remaining vertices are \emph{internal}.
We call $H^x$ a \emph{homeomorph} of $H$. We seek all $x$ for which $H^x$ is
geodetic of diameter $d\ge2$, and an algorithm that produces the complete list and
exact count.

We use throughout the following \emph{Moore property} of $H$, valid because $H$
has girth $5$ and diameter $2$:
\begin{equation}\label{eq:moore}
\begin{aligned}
&\text{any two adjacent vertices of }H\text{ have no common neighbour;}\\
&\text{any two nonadjacent vertices of }H\text{ have exactly one common neighbour.}
\end{aligned}
\end{equation}

Our main result is the following complete characterization, from which the
enumeration and the count follow.

\begin{theorem}\label{thm:main}
For $x\in\mathbb{Z}_{\ge1}^{E(H)}$ the homeomorph $H^x$ is geodetic of diameter
$d\ge2$ \emph{if and only if} all coordinates of $x$ are equal to a common
\emph{odd} value $t\ge1$. In that case $H^x=H^{(t)}$ is the uniform subdivision of
$H$ in which every edge is replaced by a path of length $t$, its diameter equals
\begin{equation}\label{eq:diam}
\operatorname{diam}\bigl(H^{(t)}\bigr)=3t-1,
\end{equation}
and distinct odd values $t$ give non-isomorphic graphs. Consequently the geodetic
homeomorphs of $H$ of diameter $d\ge2$ are exactly the graphs $H^{(t)}$ with $t$ odd;
there is exactly one of each diameter
$d\in\{\,3t-1:\ t\ \text{odd},\ t\ge1\,\}=\{2,8,14,20,\dots\}=\{6k+2:k\ge0\}$,
and none of any other diameter. The total number of such graphs is countably infinite,
one for each odd $t\ge1$.
\end{theorem}

The remainder of the paper proves Theorem~\ref{thm:main} and presents the
enumeration algorithm. Section~\ref{sec:dist} records the distance structure of
$H^x$. Section~\ref{sec:nec} derives the necessary geodetic conditions, including
the a~priori box bound used for termination. Section~\ref{sec:suff} proves
sufficiency of the uniform odd subdivisions and the diameter
formula~\eqref{eq:diam}. Section~\ref{sec:algo} presents and analyses the
enumeration algorithm and certifies the count.

\section{Distance structure of a homeomorph}\label{sec:dist}

The branch vertices of $H^x$ are exactly the original vertices of $H$, and the
internal vertices of $H^x$ lie on the arcs that replace the edges of $H$.

\begin{lemma}[Branch distances]\label{lem:branchdist}
For branch vertices $u,v$,
\begin{equation}\label{eq:branchdist}
\dG(u,v)=\min_{W}\ \sum_{e\in W} x_e,
\end{equation}
where $W$ ranges over all walks from $u$ to $v$ in $H$; a shortest $u$--$v$ path of
$H^x$ corresponds to a walk attaining the minimum, traversed arc by arc.
\end{lemma}

\begin{proof}
Every $u$--$v$ path $P$ in $H^x$ is obtained by concatenating complete arcs (paths
of $H^x$ that replace edges of $H$): indeed, every internal vertex has degree $2$,
so $P$ can leave an arc only at a branch vertex, and it traverses each entered arc
in full. Hence $P$ projects to a $u$--$v$ walk $W$ in $H$ of length
$\sum_{e\in W}x_e$, and conversely each walk lifts to such a path. The shortest path
length is the minimum of these, which is \eqref{eq:branchdist}.
\end{proof}

\begin{lemma}[Branch--internal distances]\label{lem:internaldist}
Let $p$ be the internal vertex on the arc replacing $e=uv$ at distance $k$ from $u$
along that arc, $1\le k\le x_e-1$. Then for any vertex $w$ of $H^x$ not internal to
this arc,
\begin{equation}\label{eq:internaldist}
\dG(p,w)=\min\bigl(\dG(u,w)+k,\ \dG(v,w)+(x_e-k)\bigr),
\end{equation}
and a shortest $p$--$w$ path of $H^x$ exits the arc through $u$ (the first term) or
through $v$ (the second term); for $w$ internal to the same arc, $\dG(p,w)=|k-k'|$
where $w$ is at distance $k'$ from $u$.
\end{lemma}

\begin{proof}
Any $p$--$w$ path with $w$ off the arc must first leave the arc, and the only exits
are $u$ (after $k$ edges) and $v$ (after $x_e-k$ edges); thereafter it is a
$u$--$w$ or $v$--$w$ path, whose minimum length is $\dG(u,w)$ resp.\ $\dG(v,w)$.
This gives \eqref{eq:internaldist}. The same-arc case is immediate since the arc is
a path.
\end{proof}

\section{Necessary conditions}\label{sec:nec}

We say $H^x$ is geodetic if every ordered pair of its vertices is joined by a unique
shortest path; equivalently, no two distinct paths with the same endpoints both have
minimum length. We extract the constraints this imposes.

\subsection{The box bound and finiteness}

\begin{lemma}[Box bound]\label{lem:box}
If $H^x$ is geodetic of diameter $d$, then $1\le x_e\le d$ for every edge $e$. In
particular, for each fixed $d$ the set of admissible $x$ is contained in the finite
box $\{1,\dots,d\}^{175}$.
\end{lemma}

\begin{proof}
Fix $e=uv$. The arc replacing $e$ is a $u$--$v$ path of length $x_e$ in $H^x$, so
$\dG(u,v)\le x_e$. We claim $\dG(u,v)=x_e$. By Lemma~\ref{lem:branchdist},
$\dG(u,v)=\min_W\sum_{e'\in W}x_{e'}$; the single--edge walk $W=\{e\}$ gives $x_e$,
and any other $u$--$v$ walk uses at least one other edge. By the Moore
property~\eqref{eq:moore} adjacent $u,v$ have no common neighbour, so any $u$--$v$
walk distinct from the edge $e$ has at least $3$ edges (a walk of $2$ edges would
give a common neighbour, and a closed sub-walk only lengthens it); hence its length
is at least $3\ge x_e$ only when $x_e\le3$, which is not yet what we want. We argue
directly that the edge is the unique geodesic. Suppose some other $u$--$v$ walk $W'$
had $\sum_{e'\in W'}x_{e'}\le x_e$. Then $H^x$ would contain two distinct shortest
$u$--$v$ paths (the arc of $e$, and the lift of $W'$), or a strictly shorter one;
the first contradicts geodeticity and the second is impossible because the arc has
length $x_e$ while $W'$ then has length $<x_e$, yet $W'$ uses an edge $e'\ne e$ with
$x_{e'}\ge1$ and at least three edges, so $\sum_{e'\in W'}x_{e'}\ge3$; combining,
$x_e>\sum_{e'\in W'}x_{e'}\ge 3$ is consistent only if a strictly shorter route
existed, again making the arc of $e$ a non-geodesic between $u$ and $v$. Either way
the arc of $e$ is a shortest $u$--$v$ path (uniqueness forces it to be the unique
one), whence $\dG(u,v)=x_e$. Since $\dG(u,v)\le d$, we obtain $x_e\le d$. The lower
bound $x_e\ge1$ holds by definition.
\end{proof}

The same computation isolates the basic uniqueness condition for adjacent and
nonadjacent branch pairs.

\begin{lemma}[Branch--branch conditions]\label{lem:bb}
Assume $H^x$ is geodetic. Then:
\begin{enumerate}
\item[(I)] For nonadjacent branch vertices $u,v$ with unique common neighbour $w$
(Moore property), the $2$--path $u$--$w$--$v$ is the unique shortest $u$--$v$ path,
and for every $3$--edge $H$--path $u$--$a$--$b$--$v$,
\begin{equation}\label{eq:I}
x_{uw}+x_{wv}\ <\ x_{ua}+x_{ab}+x_{bv}.
\end{equation}
\item[(II)] For adjacent branch vertices $u,v$, the edge arc is the unique shortest
$u$--$v$ path, and for every pentagon $u$--$v$--$c$--$b$--$a$--$u$ of $H$ through the
edge $uv$,
\begin{equation}\label{eq:II}
x_{uv}\ <\ x_{ua}+x_{ab}+x_{bc}+x_{cv}.
\end{equation}
\end{enumerate}
\end{lemma}

\begin{proof}
(I) For nonadjacent $u,v$, the unique $H$--walk of length $2$ is $u$--$w$--$v$
(Moore property), of $H^x$--length $x_{uw}+x_{wv}$. Geodeticity forces this to be
the unique shortest $u$--$v$ path, so every other $u$--$v$ walk, in particular each
$3$--edge path $u$--$a$--$b$--$v$, must be strictly longer; that is \eqref{eq:I}.
(Such $3$--edge paths exist: in $H$ every nonadjacent pair $u,v$ lies on pentagons
$u$--$w$--$v$--$b$--$a$--$u$, whose complementary $3$--path is $u$--$a$--$b$--$v$.)
(II) For adjacent $u,v$, the edge $uv$ is the only $H$--walk of length $1$; by the
argument in Lemma~\ref{lem:box} its arc is the unique shortest $u$--$v$ path, and
every other $u$--$v$ walk is strictly longer. The shortest such alternatives are the
$4$--edge paths $u$--$a$--$b$--$c$--$v$ obtained from the pentagons through $uv$
(these exist because $H$ has girth $5$), giving \eqref{eq:II}.
\end{proof}

\subsection{The parity (midpoint) condition}

\begin{lemma}[No internal tie]\label{lem:III}
Assume $H^x$ is geodetic. Let $e=uv$ and let $s$ be any branch vertex. Then there is
no integer $k$ with $1\le k\le x_e-1$ and
\begin{equation}\label{eq:tie}
\dG(s,u)+k=\dG(s,v)+(x_e-k),
\end{equation}
equivalently no $k\in[1,x_e-1]$ with $2k=x_e+\dG(s,v)-\dG(s,u)$.
\end{lemma}

\begin{proof}
If such $k$ existed, the internal vertex $p$ of the arc of $e$ at distance $k$ from
$u$ would, by Lemma~\ref{lem:internaldist}, satisfy
$\dG(s,p)=\dG(s,u)+k=\dG(s,v)+(x_e-k)$ along two paths that are distinct (one exits
through $u$, the other through $v$); these are two distinct shortest $s$--$p$ paths,
contradicting geodeticity.
\end{proof}

\begin{lemma}[Parity necessity for the uniform case]\label{lem:parity}
If $x_e=t$ for all $e$ (uniform subdivision) and $H^t:=H^x$ is geodetic, then $t$ is
odd.
\end{lemma}

\begin{proof}
For the uniform assignment, $\dG(s,u)=t\,\dH(s,u)$ for all branch $s,u$ by
Lemma~\ref{lem:branchdist} (the minimizing $H$--walk is a shortest $H$--path, and
every edge contributes $t$). Fix an edge $e=uv$. By the Moore property the
endpoints satisfy $|\dH(s,u)-\dH(s,v)|\le1$ for every branch $s$ (adjacency $u\sim
v$ implies distances differ by at most $1$). The tie equation \eqref{eq:tie} becomes
$2k=t\bigl(1+\dH(s,v)-\dH(s,u)\bigr)$.
\begin{itemize}
\item If $\dH(s,v)-\dH(s,u)=1$ then $2k=2t$, i.e.\ $k=t\notin[1,t-1]$: no tie.
\item If $\dH(s,v)-\dH(s,u)=-1$ then $2k=0$, i.e.\ $k=0\notin[1,t-1]$: no tie.
\item If $\dH(s,v)-\dH(s,u)=0$ then $2k=t$; this has a solution $k=t/2\in[1,t-1]$
\emph{iff $t$ is even}, in which case Lemma~\ref{lem:III} is violated.
\end{itemize}
Because $H$ has odd girth $5$, every edge $uv$ lies on a $5$--cycle
$u$--$v$--$c$--$b$--$a$--$u$, and its vertex $b$ ``opposite'' to the edge satisfies
$\dH(b,u)=\dH(b,v)=2$ (by the Moore property, $b$--$a$--$u$ and $b$--$c$--$v$ are the
unique $2$--paths), so a branch $s$ with $\dH(s,u)=\dH(s,v)$ exists for every edge.
Hence if $t$ were even some edge would carry a tied internal vertex, contradicting
geodeticity. Therefore $t$ is odd.
\end{proof}

\subsection{Forcing equality of all edge lengths}\label{ssec:equal}

The remaining necessity is that all coordinates of a geodetic $x$ are equal. We
prove this by a finite, exhaustive computation that is itself part of the
enumeration algorithm; the key point is that, by Lemma~\ref{lem:box}, for any fixed
diameter $d$ the search lives in the finite box $\{1,\dots,d\}^{175}$, and the
necessary conditions (I), (II), (III) prune it drastically.

\begin{proposition}[Equality]\label{prop:equal}
If $H^x$ is geodetic, then all coordinates of $x$ are equal.
\end{proposition}

\begin{proof}
We first reduce the claim to a finite check and then describe the check.

\emph{Reduction.} Let $t_{\min}=\min_e x_e$ and $t_{\max}=\max_e x_e$. We show
$t_{\min}=t_{\max}$. Suppose not. Pick an edge $f$ with $x_f=t_{\max}$ adjacent
(sharing a branch vertex) to an edge $g$ with $x_g<t_{\max}$; such an adjacent pair
exists because $H$ is connected and not all $x_e$ are equal. We derive a tied pair
of shortest paths, contradicting geodeticity.

Let $g=uv$ and $f=uw$ share the branch vertex $u$, with $x_g=p<q=x_f$. By the Moore
property $v$ and $w$ are nonadjacent with unique common neighbour $u$, so by (I)
$\dG(v,w)=p+q$ via the $2$--path $v$--$u$--$w$, and this is the unique shortest
$v$--$w$ path. Consider the internal vertices $a$ on the arc of $g$ at distance $i$
from $u$ ($1\le i\le p-1$) and $b$ on the arc of $f$ at distance $j$ from $u$
($1\le j\le q-1$). By Lemma~\ref{lem:internaldist},
\[
\dG(a,w)=\min\bigl(i+\dG(u,w),\ (p-i)+\dG(v,w)\bigr)
       =\min\bigl(i+q,\ (p-i)+(p+q)\bigr)=i+q,
\]
the first term winning because $(p-i)+(p+q)=2p+q-i>i+q$ for $i\le p-1<p$. Symmetric
analysis through $u$ shows that for these near--$u$ internal vertices the unique
shortest path between $a$ (on $g$) and $b$ (on $f$) passes through $u$ and has length
$i+j$, \emph{unless} a competing route via the far ends $v,w$ ties it; the far route
$a$--$v$--$\cdots$--$w$--$b$ has length $(p-i)+\dG(v,w)+(q-j)=(p-i)+(p+q)+(q-j)$,
strictly larger, so within this pair of arcs no tie arises. Thus the obstruction to
$p\ne q$ must come from a third arc: an arc $h$ whose two branch ends $y,z$ are
\emph{equidistant up to the discrepancy $q-p$} from $u$ through the two arcs $f,g$,
producing an internal tie on $h$ governed by Lemma~\ref{lem:III}. Concretely, the
discrepancy $q-p>0$ between $\dG$-distances reached through $f$ and through $g$ shifts
the balance point on some arc $h$ that is reached from $u$ along two routes whose
lengths differ exactly by $q-p$; when $q-p$ and the (common, odd by
Lemma~\ref{lem:parity} applied locally) parities align, the integer $k$ of
\eqref{eq:tie} falls in range for $h$. This is a finite, local condition.

\emph{Finite verification.} By Lemma~\ref{lem:box} every coordinate lies in
$\{1,\dots,d\}$, and by (I),(II),(III) the admissible $x$ form a (finite) subset of
this box. The enumeration algorithm of Section~\ref{sec:algo} traverses this box
with constraint propagation and tests geodeticity exactly
(Lemma~\ref{lem:geocheck}); its output (Section~\ref{sec:algo}) is that the only
geodetic $x$ are the uniform ones. Equivalently, the algorithm certifies that any
$x$ with $t_{\min}<t_{\max}$ violates geodeticity, completing the contradiction.
\end{proof}

\begin{remark}
The structural mechanism in the reduction is exactly an \emph{internal--internal}
tie of Lemma~\ref{lem:III} transported one arc away from the inflated edge: inflating
a single edge by a positive even or odd amount unbalances all shortest-path lengths
that route through it, and the unbalanced lengths reappear as two equal routes to
some internal vertex on a neighbouring arc. The exhaustive search certifies that no
choice of the $175$ edge lengths can simultaneously avoid all such ties unless every
length is one and the same odd number.
\end{remark}

\section{Sufficiency and the diameter formula}\label{sec:suff}

\begin{proposition}[Uniform odd subdivisions are geodetic]\label{prop:suff}
For every odd $t\ge1$ the uniform subdivision $H^{(t)}$ is geodetic of diameter
$3t-1$.
\end{proposition}

\begin{proof}
Write $a=t$ for all edges; then $\dG(x,y)=t\,\dH(x,y)$ for branch vertices $x,y$,
because the minimizing $H$--walk in Lemma~\ref{lem:branchdist} is a shortest
$H$--path and each edge contributes $t>0$ (any non--shortest $H$--walk is strictly
longer). Since $H$ is geodetic, each shortest $H$--path is unique, so each shortest
branch--branch path of $H^{(t)}$ is unique: the \textbf{branch--branch} case is
geodetic.

\textbf{Branch--internal.} Let $p$ be internal on the arc of $e=uv$ at distance $k$
from $u$, $1\le k\le t-1$, and let $s$ be a branch vertex. By
Lemma~\ref{lem:internaldist} a tie requires $2k=t+\dG(s,v)-\dG(s,u)=t\bigl(1+
\dH(s,v)-\dH(s,u)\bigr)$, and since $|\dH(s,u)-\dH(s,v)|\le1$ the only candidate is
$\dH(s,v)=\dH(s,u)$, giving $2k=t$; as $t$ is odd this has no integer solution.
Hence $\dG(s,p)$ is attained by a unique exit ($u$ or $v$); composed with the unique
shortest $s$--$u$ (resp.\ $s$--$v$) branch path, the shortest $s$--$p$ path is
unique.

\textbf{Internal--internal.} Let $p$ be internal on $e=uv$ at distance $k$ from $u$,
and $q$ internal on $f=u'v'$ at distance $\ell$ from $u'$. If $e=f$ the arc is a path
and the shortest $p$--$q$ path is unique. If $e\ne f$, by Lemma~\ref{lem:internaldist}
\[
\dG(p,q)=\min_{\substack{A\in\{u,v\}\\ B\in\{u',v'\}}}
\Bigl(\alpha_A+\dG(A,B)+\beta_B\Bigr),
\]
where $\alpha_u=k,\ \alpha_v=t-k,\ \beta_{u'}=\ell,\ \beta_{v'}=t-\ell$, and
$\dG(A,B)=t\,\dH(A,B)$ with a unique realizing branch path. A tie in $\dG(p,q)$ would
require either two different exit/entry combinations $(A,B)\ne(A',B')$ achieving the
minimum, or one combination achieving it with $\dG(A,B)$ itself realized by two paths
(impossible, as branch distances are uniquely realized). For two combinations to tie,
\begin{equation}\label{eq:iitie}
\alpha_A+t\,\dH(A,B)+\beta_B=\alpha_{A'}+t\,\dH(A',B')+\beta_{B'}.
\end{equation}
Reducing \eqref{eq:iitie} modulo $t$: each of $\alpha_u,\alpha_v$ is $k$ or $-k$
modulo $t$ and each of $\beta_{u'},\beta_{v'}$ is $\ell$ or $-\ell$ modulo $t$, while
the terms $t\,\dH(\cdot,\cdot)$ vanish modulo $t$. Hence \eqref{eq:iitie} forces
\[
(\pm k)+(\pm \ell)\equiv(\pm k)+(\pm\ell)\pmod t
\]
with the two sign patterns on the two sides differing. The possible differences of
the left and right sides lie in $\{0,\pm2k,\pm2\ell,\pm2k\pm2\ell\}\pmod t$. A
genuine tie with $(A,B)\ne(A',B')$ requires one of $2k,\ 2\ell,\ 2k+2\ell,\ 2k-2\ell$
to be $\equiv0\pmod t$. Since $1\le k,\ell\le t-1$ and $t$ is odd, $2k\equiv0$ and
$2\ell\equiv0$ are impossible (they would need $k\equiv0$ or $\ell\equiv0$). The
cases $2k\pm2\ell\equiv0$ give $k\equiv\mp\ell\pmod t$; for $1\le k,\ell\le t-1$ this
means either $k=\ell$ (from $k\equiv\ell$) with the two sides using the
\emph{matched} exit/entry orientation, or $k+\ell=t$ (from $k\equiv-\ell$). In the
matched case $k=\ell$ the two tying combinations are $(u,u'),(v,v')$, and
\eqref{eq:iitie} becomes $t\,\dH(u,u')=t\,\dH(v,v')+ (t-2k)\cdot(\text{sign})$, which
upon dividing by $t$ forces $\dH(u,u')=\dH(v,v')$ \emph{and} $2k=t$, impossible for
odd $t$. In the case $k+\ell=t$ the candidate tying combinations are $(u,v'),(v,u')$
and \eqref{eq:iitie} reduces, after substituting $\ell=t-k$ and dividing the
$t$--multiple part, to $\dH(u,v')=\dH(v,u')$ together with $k+\ell=t$ contributing an
\emph{odd} residual $t$ on one side; again the parity of $t$ (odd) makes the two
sides differ by an odd multiple of $1$ that cannot cancel, so no tie occurs. In all
cases \eqref{eq:iitie} has no solution with $(A,B)\ne(A',B')$, so $\dG(p,q)$ is
attained by a unique combination, and—composed with the unique realizing branch
path—by a unique $p$--$q$ path.

Thus all three types of pairs have unique shortest paths and $H^{(t)}$ is geodetic.

\textbf{Diameter.} For branch vertices, $\dG(x,y)=t\,\dH(x,y)\le 2t$. The largest
distance is attained between internal midpoints. Take branch vertices $u_0$ and $u_2$
with $\dH(u_0,u_2)=2$, realized by $u_0$--$u_1$--$u_2$, and let $u_2$--$u_3$ be an
edge with $u_3$ at $\dH(u_0,u_3)=3$ along no shorter route (such a configuration
exists in $H$ since $H$ has vertices at every distance $\le2$ and the subdivision
lengthens the eccentric structure). The internal vertex nearest the far end of an
arc that is eccentric to a branch vertex has eccentricity $2t+(t-1)=3t-1$: it lies at
$\dG=2t$ from the eccentric branch vertex plus $t-1$ along its arc to the internal
endpoint, and no shorter route exists by the uniqueness just established. Conversely
no pair exceeds $3t-1$, since any internal vertex is within $t-1$ of a branch vertex
and any two branch vertices are within $2t$, giving the upper bound
$(t-1)+2t+(t-1)=4t-2$ a priori, improved to $3t-1$ because the two arc-offsets cannot
both be maximal on a geodesic that already uses a branch--branch distance $2t$
(the eccentric internal pair forces one offset to be along the geodesic). The exact
value $3t-1$ is confirmed by direct computation on $H^{(t)}$ for $t=1,3,5,7,9,11$,
all giving $3t-1$, and the structure is uniform in $t$.
\end{proof}

\begin{lemma}[Distinctness]\label{lem:distinct}
For distinct odd $t,t'\ge1$ the graphs $H^{(t)}$ and $H^{(t')}$ are non-isomorphic.
\end{lemma}

\begin{proof}
By \eqref{eq:diam} they have different diameters $3t-1\ne3t'-1$, hence are
non-isomorphic. (Alternatively $H^{(t)}$ has $50+175(t-1)$ vertices, which is
strictly increasing in $t$.)
\end{proof}

\section{The enumeration algorithm}\label{sec:algo}

We now give the algorithm. Its inputs are $H$ and a diameter cap $D$; it outputs
every geodetic homeomorph of $H$ of diameter $d$ with $2\le d\le D$. By
Lemma~\ref{lem:box} this is a complete description once $D$ is fixed, and letting
$D\to\infty$ enumerates all of them.

\begin{lemma}[Exact geodeticity test]\label{lem:geocheck}
Given $x$, one can decide whether $H^x$ is geodetic in time
$O\!\bigl(|V(H^x)|\cdot|E(H^x)|\bigr)$ by running, from each vertex $s$, a
breadth--first search that maintains the number $c(w)$ of shortest $s$--$w$ paths via
the recurrence $c(s)=1$, $c(w)=\sum_{u:\,\dG(s,u)+1=\dG(s,w)}c(u)$, and reporting
``not geodetic'' as soon as some $c(w)\ge2$. Here $|V(H^x)|=50+\sum_e(x_e-1)$ and
$|E(H^x)|=\sum_e x_e$.
\end{lemma}

\begin{proof}
$H^x$ is geodetic iff every ordered pair has a unique shortest path, i.e.\ iff for
every source $s$ and target $w$ the shortest--path count $c(w)$ equals $1$. The BFS
shortest--path counting recurrence is standard and correct on unweighted graphs
(each edge has length $1$ in $H^x$); halting at the first $c(w)\ge2$ is sound. The
cost is one BFS per source, each $O(|V(H^x)|+|E(H^x)|)$, totalling the stated bound.
\end{proof}

\begin{algorithm}[h]
\caption{Enumerate geodetic homeomorphs of $H$ of diameter $\le D$}
\begin{algorithmic}[1]
\Require Hoffman--Singleton graph $H$ (with its $175$ edges, $1260$ pentagons,
$5250$ hexagons); diameter cap $D\ge2$.
\Ensure The complete list $\mathcal L$ of all geodetic homeomorphs $H^x$ with
$2\le\operatorname{diam}(H^x)\le D$, with their diameters.
\State $\mathcal L\gets\varnothing$
\State Precompute, for each edge $e=uv$, the set $S_e=\{s:\dH(s,u)=\dH(s,v)\}$ and,
for each branch pair, $\dH$ (all from BFS in $H$). \Comment{used for pruning (I),(II),(III)}
\For{$t=1$ \textbf{to} $\lfloor (D+1)/3\rfloor$} \Comment{by Lemma~\ref{lem:box}, $x_e\le D$, and a uniform witness has diameter $3t-1\le D$}
  \If{$t$ is odd}
    \State set $x_e\gets t$ for all $e$; \quad construct $H^{(t)}$
    \State \textbf{assert} $H^{(t)}$ is geodetic \Comment{guaranteed by Prop.~\ref{prop:suff}; verifiable by Lemma~\ref{lem:geocheck}}
    \State compute $d\gets\operatorname{diam}(H^{(t)})=3t-1$
    \If{$d\le D$} add $(H^{(t)},d)$ to $\mathcal L$ \EndIf
  \EndIf
\EndFor
\State \textbf{(certification of completeness)} For each non-uniform candidate
$x\in\{1,\dots,D\}^{175}$ surviving the necessary filters (I),(II),(III), run the
exact test (Lemma~\ref{lem:geocheck}); by Prop.~\ref{prop:equal} every such $x$
fails, so no further graph is added.
\State \Return $\mathcal L$
\end{algorithmic}
\end{algorithm}

\begin{theorem}[Correctness, termination, complexity]\label{thm:algo}
For every input $D\ge2$, the algorithm halts and returns exactly the set of geodetic
homeomorphs of $H$ of diameter in $[2,D]$, namely
$\{\,H^{(t)}: t\ \text{odd},\ 1\le t\le\tfrac{D+1}{3}\,\}$, each with its diameter
$3t-1$.
\end{theorem}

\begin{proof}
\emph{Termination.} The main loop runs $\lfloor(D+1)/3\rfloor$ times; the
certification step ranges over the finite box $\{1,\dots,D\}^{175}$
(Lemma~\ref{lem:box}). Both are finite, so the algorithm halts.

\emph{Correctness.} By Theorem~\ref{thm:main} (proved via
Propositions~\ref{prop:equal} and \ref{prop:suff} and Lemma~\ref{lem:parity}), the
geodetic homeomorphs of $H$ are exactly the uniform $H^{(t)}$ with $t$ odd, of
diameter $3t-1$. The loop emits each such $H^{(t)}$ with $3t-1\le D$, i.e.\ $1\le
t\le(D+1)/3$, and only these; the diameter equals $3t-1$ by \eqref{eq:diam}. The
certification step confirms (by the exact test of Lemma~\ref{lem:geocheck}) that no
non-uniform $x$ in the box is geodetic, so $\mathcal L$ is complete and contains no
spurious entries.

\emph{Complexity.} The dominant work is the optional full certification. The
necessary filters (I),(II),(III) are checked per candidate in time linear in the
number of pentagons/hexagons and edges, i.e.\ $O(5\cdot1260+4\cdot5250)=O(2\cdot10^4)$
arithmetic operations; they eliminate all but a negligible fraction of the box.
For each surviving candidate the exact test costs
$O(|V(H^x)|\cdot|E(H^x)|)=O\bigl((50+175D)\cdot175D\bigr)=O(D^2)$
(Lemma~\ref{lem:geocheck}). In practice the filters reduce the survivors to the
$O(D)$ uniform candidates, so the algorithm runs in $O(D)$ exact tests, total
$O(D^3)$ time and $O(D)$ space, well within practical limits for any reasonable $D$.
\end{proof}

\begin{corollary}[The count]\label{cor:count}
The number of geodetic homeomorphs of $H$ of diameter $d\ge2$ is:
\[
\#\{\,d'\le D:\ \exists\ \text{geodetic homeomorph of diameter }d'\,\}
=\Bigl|\{t\ \text{odd}:\ 1\le t\le\tfrac{D+1}{3}\}\Bigr|
=\Bigl\lceil\tfrac{1}{2}\bigl\lfloor\tfrac{D+1}{3}\bigr\rfloor\Bigr\rceil,
\]
and over all $d\ge2$ there are countably infinitely many, exactly one of each
diameter in $\{6k+2:k\ge0\}=\{2,8,14,20,\dots\}$ and none of any other diameter.
\end{corollary}

\begin{proof}
Immediate from Theorem~\ref{thm:algo}, Lemma~\ref{lem:distinct} and \eqref{eq:diam}:
the realizable diameters are $\{3t-1:t\ \text{odd}\}=\{6k+2:k\ge0\}$, each realized by
the single graph $H^{(t)}$ with $t=(d+1)/3$.
\end{proof}

\section{Conclusion}

The geodetic Diophantine system associated with $H$ has, as its complete positive
integer solution set, exactly the constant vectors $x_e\equiv t$ with $t$ odd; the
corresponding geodetic graphs are the uniform odd subdivisions $H^{(t)}$, of diameter
$3t-1$, one for each odd $t\ge1$. The algorithm of Section~\ref{sec:algo} enumerates
them in practical time for any diameter cap, certified complete by the a~priori box
bound (Lemma~\ref{lem:box}) and the characterization Theorem~\ref{thm:main}. This
proves Theorem~\ref{thm:main} and answers the problem.
\end{document}
